%%==========================================================================%%
%% methods.tex -- "Methods". Subsection structure mirrors the AutoPET         %%
%% descriptor (Gatidis 2022), adapted to a pediatric whole-body FDG-PET/MRI   %%
%% lymphoma lesion dataset. American spelling. [TBD] = our specifics, mostly  %%
%% from the DICOM headers. No results or analyses anywhere in the descriptor.      %%
%%==========================================================================%%

\section*{Methods}\label{sec:methods}

\subsection*{Ethics statement}\label{subsec:ethics}
The collection and retrospective use of these data were approved by the Stanford University Institutional Review Board under protocol 44706 (``Pediatric PET/MR Image Registry''), with approval granted on [TBD: approval date], and all procedures were carried out in accordance with the Declaration of Helsinki and the Health Insurance Portability and Accountability Act (HIPAA).
The registry operates under an approved waiver of authorization and a waiver of assent for minors under 45~CFR~46.404 [TBD: confirm the exact waiver wording and the public-release modification].
The imaging was de-identified to a HIPAA standard before release.
The DICOM headers were de-identified using an automated pipeline [TBD: tool and version], and the pixel data were processed to remove identifiable features, including defacing of the face and of the pelvic region on every spatially co-registered series, because the co-registration of PET and MRI can otherwise re-expose a region defaced on one modality alone [TBD: defacing tool and version]~\cite{Selfridge_2023}.
The standard applied for release was a HIPAA Expert Determination [TBD: certifying party], under which the dataset was certified to carry no more than a very small risk of re-identification.
The resulting access conditions are described in the Usage Notes and the Data Availability statement.

\subsection*{Participants}\label{subsec:cohort}
Subjects were retrospectively identified at Lucile Packard Children's Hospital at Stanford, a single center.
The inclusion criterion was a clinically indicated whole-body FDG-PET/MRI examination performed for pediatric or young-adult lymphoma (Hodgkin or non-Hodgkin), in patients aged [TBD] years (median [TBD]).
Examinations were acquired between [TBD: start] and [TBD: end].
The sex distribution was [TBD: N female, N male].
[TBD: formal exclusion criteria with counts, for example incomplete whole-body coverage, non-diagnostic image quality, or PET and MRI series that could not be spatially aligned.]
[TBD: number of Hodgkin versus non-Hodgkin cases.]

The released examinations comprise two labeled subsets, which may overlap at the patient level, so a patient may contribute to both and the subset counts need not sum to the cohort total.
The first is a lesion-positive subset of examinations with manually drawn FDG-avid tumor-lesion segmentations, comprising approximately [TBD] examinations from [TBD] patients.
The second is a disease-free reference subset of follow-up examinations, each demonstrating complete metabolic response with no FDG-avid malignant lesions, comprising approximately [TBD] examinations.
The disease-free reference subset is the pathology-negative analog of the negative-control studies in the adult whole-body FDG-PET/CT lesion dataset~\cite{Gatidis_2022}.
A subject-by-subject description of demographic and clinical characteristics is provided in the accompanying metadata (Section~\ref{sec:datarecords}), and patient characteristics are summarized in Table~\ref{tab:cohort}.

%%-- Patient-characteristics table placeholder (cf. AutoPET Table 1) --------%%
\begin{table}[h]
\caption{Patient characteristics of the released cohort, by lymphoma class and labeled subset. All cells are [TBD] pending reconciliation of the release manifest.}\label{tab:cohort}
\begin{tabular}{@{}lllll@{}}
\toprule
 & \multicolumn{2}{c}{Lesion-positive subset} & \multicolumn{2}{c}{Disease-free reference subset} \\
\cmidrule(lr){2-3}\cmidrule(lr){4-5}
Characteristic & Hodgkin & Non-Hodgkin & Hodgkin & Non-Hodgkin \\
\midrule
Patients ($n$)            & [TBD] & [TBD] & [TBD] & [TBD] \\
Examinations ($n$)        & [TBD] & [TBD] & [TBD] & [TBD] \\
Female / male             & [TBD] & [TBD] & [TBD] & [TBD] \\
Age, years (mean $\pm$ SD) & [TBD] & [TBD] & [TBD] & [TBD] \\
\botrule
\end{tabular}
\end{table}

\subsection*{Image acquisition}\label{subsec:acquisition}
Examinations were acquired as integrated whole-body FDG-PET/MRI on a [TBD: PET/MR scanner manufacturer and model] operating at a field strength of [TBD: field strength].

For the PET component, [$^{18}$F]fluorodeoxyglucose was injected intravenously in all examinations, at a weight-based pediatric activity of [TBD: injected activity, following weight-based reference levels]~\cite{Lassmann_2007}.
PET acquisition was initiated after an uptake period of [TBD: uptake time].
PET images were reconstructed using [TBD: reconstruction algorithm and parameters, for example iterations, subsets, filter, matrix size, and voxel size].
Attenuation correction was MR-based, derived from a dedicated MRI sequence rather than from CT [TBD: attenuation-correction method].
The implications of MR-based attenuation correction for quantitative uptake measures are noted in the Usage Notes.

For the MRI component, because protocols were determined by clinical indication, the number and type of sequences vary across subjects and include both non-contrast and contrast-enhanced acquisitions.
Whole-body coverage was obtained using [TBD: acquisition strategy, for example multi-station composed sequences], and the acquired contrasts include [TBD: for example T1-weighted, T2-weighted, STIR, Dixon, and diffusion-weighted imaging].
Acquisition parameters for the most frequently applied sequences are summarized in Table~\ref{tab:acq}, and the complete per-examination parameters are recorded in the DICOM headers.

%%-- Acquisition-parameter table placeholder (cf. AutoPET acquisition) ------%%
\begin{table}[h]
\caption{Representative PET and MRI acquisition parameters for the most frequently applied protocol. All values are [TBD] pending extraction from the DICOM headers.}\label{tab:acq}
\begin{tabular}{@{}llll@{}}
\toprule
Parameter & PET & [MRI seq.\ 1] & [MRI seq.\ 2] \\
\midrule
Tracer / activity        & FDG, [TBD] & n/a & n/a \\
Uptake time [min]        & [TBD] & n/a & n/a \\
Reconstruction           & [TBD] & n/a & n/a \\
Attenuation correction   & MR-based [TBD] & n/a & n/a \\
TR [ms]                  & n/a & [TBD] & [TBD] \\
TE [ms]                  & n/a & [TBD] & [TBD] \\
Slice thickness [mm]     & [TBD] & [TBD] & [TBD] \\
Voxel size [mm]          & [TBD] & [TBD] & [TBD] \\
\botrule
\end{tabular}
\end{table}

\subsection*{Data processing}\label{subsec:processing}
The released imaging is provided in its native DICOM format, with the acquisition parameters retained in the DICOM headers.
Conversion utilities are provided that read the DICOM imaging and produce derived volumes (NIfTI), apply the standardized-uptake-value (SUV) conversion to the PET using the activity, body weight, and timing fields recorded in the headers, and resample the PET and MRI series onto a common grid where required (Section~\ref{sec:code}).
The manual lesion and disease-free-reference masks were converted from their annotation format (NIfTI/NRRD) to DICOM-SEG using dcmqi [TBD: version] for release.

\subsection*{Annotation}\label{subsec:annotation}
Tumor-lesion segmentations were drawn manually in 3D~Slicer [TBD: version].
Candidate examinations were identified from the clinical record, and all FDG-avid lesions reported as malignant in the clinical PET/MRI report were segmented on the FDG-avid PET, in a slice-by-slice manner [TBD: confirm the segmentation reference series against the DICOM headers], a procedure that parallels the manual PET segmentation of the adult whole-body FDG-PET/CT lesion dataset~\cite{Gatidis_2022}.
Annotation was performed by a team of trained readers [TBD: number and experience of annotators, and the annotation standard operating procedure].
The disease-free reference examinations carry no lesion masks by definition, and they are recorded as complete-metabolic-response follow-up studies so that they can be selected as pathology-negative references.
The released labels record, for each examination, the target structure and the label type (voxel-level segmentation).

\subsection*{Data Properties}\label{subsec:properties}
This subsection summarizes the lesion-burden distributions of the lesion-positive subset and is the only part of this descriptor that reports summary statistics.
The Data Records carry file-level information only.
For the lesion-positive subset, the metabolic tumor volume (MTV), mean and maximum SUV, and total lesion glycolysis (TLG) of the segmented lesions amounted to (mean~$\pm$~SD) [TBD], [TBD], and [TBD], respectively [TBD: report by lymphoma class where the subset sizes permit].
SUV measurements were referenced to the liver and the mediastinal blood pool, the regions used for response assessment of FDG-avid lymphoma~\cite{Barrington_2017}.
These distributions are shown in Figure~\ref{fig:cohort}.
